Table~\ref{tab:ablation} reports three ablations. Each variant was trained with the same budget of 200 steps as a reference SALS model and evaluated on the held-out seasons.

\begin{itemize}
\item \emph{Without stress features}: the control network sees only phenology, time, design and crop, with soil water, ARID and the maximum-temperature, irradiance and ET$_0$ forecasts masked. At similar mean retention, this variant produced 28\% less electricity (700 vs 971\,MWh\,ha$^{-1}$\,yr$^{-1}$) and lowered the LER from 1.59 to 1.40 at present climate and from 1.67 to 1.42 at +3\,\textdegree C. A stress-blind policy performs like the expert rules: it learns \emph{when} in the season to share light, but it cannot respond to the heat and drought of a particular season.
\item \emph{Without the design network}: every site uses a GCR of 0.3, close to the mean GCR chosen by SALS, and only the controller is learned. The LER fell to 1.42 and 1.43, and compliance dropped to 31--38\%. The benefit of SALS therefore comes from matching the array density to the site's climate and crop: denser at hot, dry sites where shade is cheap, sparser where light limits growth. The same mean density everywhere does not achieve this.
\item \emph{Trained on the baseline climate only}: the networks were trained on the present climate alone and evaluated in all four climates. They did not learn to exploit the heat-protection value of shade. Their electricity stayed at about 670\,MWh\,ha$^{-1}$\,yr$^{-1}$ and their LER at 1.38--1.40 under all warming levels, whereas the multi-climate model's LER rose from 1.59 to 1.67. This variant was weaker even under the present climate. With a single training seed we cannot separate a regularising effect of the more diverse multi-climate data from optimisation variance, but the result shows that training on projected climates is needed for a policy that remains effective as the climate warms.
\end{itemize}

Training longer (400 instead of 200 steps) added a further 0.04--0.08 to the LER.

\begin{table}[t]
\centering
\caption{Ablation on held-out seasons. $E$: electricity (MWh\,ha$^{-1}$\,yr$^{-1}$); $r$: mean yield retention; C: compliance (\% of seasons with $r\ge0.9$). All ablations were trained for 200 steps and should be compared with SALS (200 steps).}
\label{tab:ablation}
\small
\setlength{\tabcolsep}{4pt}
\begin{tabular}{lcccccccc}
\toprule
 & \multicolumn{4}{c}{Baseline climate} & \multicolumn{4}{c}{+3\,\textdegree C}\\
\cmidrule(lr){2-5}\cmidrule(lr){6-9}
Variant & $E$ & $r$ & C (\%) & LER & $E$ & $r$ & C (\%) & LER\\
\midrule
SALS, full training (400 steps) & 1019 & 0.909 & 49.5 & 1.63 & 1173 & 0.901 & 45.1 & 1.75\\
\midrule
SALS (200 steps) & 971 & 0.897 & 47.3 & 1.59 & 1079 & 0.892 & 43.5 & 1.67\\
w/o stress features & 700 & 0.898 & 33.4 & 1.40 & 706 & 0.906 & 42.7 & 1.42\\
w/o design network (GCR 0.3) & 734 & 0.889 & 31.2 & 1.42 & 719 & 0.900 & 38.1 & 1.43\\
trained on baseline climate only & 668 & 0.904 & 41.4 & 1.38 & 676 & 0.910 & 46.6 & 1.40\\
\bottomrule
\end{tabular}
\end{table}

Finally, we retrained SALS (200 steps) with a softer floor of $\rho=0.8$ and re-evaluated all baselines with the same floor. The baselines' GCR grids and rules were recalibrated for $\rho=0.8$; the static rule then selects GCR 0.15--0.30. Under the present climate, the static design reached an LER of 1.35--1.36, seasonal sharing 1.59 and the stress rule 1.68 (1184\,MWh\,ha$^{-1}$\,yr$^{-1}$). SALS chose much denser arrays (mean GCR 0.55) and produced 1604\,MWh\,ha$^{-1}$\,yr$^{-1}$, 15\% more than a conventional PV plant with GCR 0.4, at a mean retention of 0.815. Its LER was 1.97, rising to 2.03 at +3\,\textdegree C. The ranking of the methods is therefore robust to the choice of the food-security floor. So is the main weakness of SALS: only 38--41\% of its seasons met the 80\% floor, against 74--85\% for the rules. The floor thus steers where SALS operates on the energy--food frontier, but a penalised average does not guarantee it season by season.
